\chapter{Outer Loop: Linear-Acceleration Control}\label{ch:outer_loop}

This chapter describes the outer loop.
It turns a position, velocity or acceleration setpoint into an attitude setpoint and a collective thrust.
The inner loop of \cref{ch:inner_loop} then follows the attitude.
The chapter goes from the problem, to the idea, to the design, and ends with the parameters a user can set.
The margins of the finished loop are analysed in \cref{ch:analysis}.

\section{Problem and idea}\label{sec:outer_problem}

The stock PX4 position controller turns a commanded acceleration into a thrust vector.
It uses a model of the thrust.
The model says how much thrust a given motor command gives.
It is tuned for the hover thrust of the airframe.
If the battery sags or a motor weakens, the real thrust is lower than the model says.
The vehicle then accelerates less than commanded.
The stock controller closes this gap with an integrator in the velocity loop.
The integrator is slow, and it can wind up.

The INDI law of \cref{sec:indi_principle} works in a different way.
It does not need the thrust to match a model.
It measures the acceleration that the vehicle has now.
It adds to the present thrust the thrust that would remove the acceleration error.
Four questions follow.
\begin{enumerate}
  \item Where does the commanded acceleration $\aref$ come from?
  \item What is the thrust that is acting now, and how is it measured?
  \item How do the measured acceleration and the thrust stay in step?
  \item How is a standing thrust gap removed without an integrator?
\end{enumerate}
The sections below answer them in this order.
Later sections cover the measured and modelled state, the flight modes, the ground and take-off ramp, the estimator resets and the guards.

\paragraph{Placement.}
The outer-loop module runs instead of the stock position controller when \param{MC_INDI_POS_EN} is 1 (reboot required).
Only one of the two runs.
The stock position controller still runs inside the module, but only to produce the acceleration reference $\aref$.
Its thrust output is never published.
The module reuses the stock setpoint handling, limits, take-off ramp and goto logic.
Only the step from acceleration to thrust is new.
It serves every multicopter mode in which the stock position controller is active (\cref{sec:outer_modes}).
The attitude loop, the rate loop and the control allocator downstream keep their roles.
The rate loop receives one extra signal, a rate feed-forward (\cref{sec:outer_ff}).

\Cref{tab:outer_io} lists the signals.

\begin{table}[tbp]
\caption{Signals of the outer-loop module.}\label{tab:outer_io}
\centering
\small
\begin{tabularx}{\textwidth}{@{}L{4.4cm} Y@{}}
\toprule
Signal (uORB topic) & Use \\
\midrule
\topic{trajectory_setpoint}, \topic{goto_setpoint}, \topic{vehicle_constraints} & Input: position, velocity, acceleration and yaw references. \\
\topic{vehicle_local_position}, \topic{hover_thrust_estimate} & Input: position and velocity from the estimator, and the thrust model scale. \\
\topic{vehicle_attitude} & Input: rotation $R$ and yaw. \\
\topic{vehicle_acceleration} & Input: measured specific force, handled on every sample (for example about 267~Hz in the reference simulation). \\
\topic{esc_status} & Input: rotor speed, for the measured thrust. \\
\topic{vehicle_thrust_setpoint} & Input: collective used to restart the state. \\
\topic{control_allocator_status} & Input: saturation flag. \\
\topic{vehicle_control_mode}, \topic{vehicle_land_detected} & Input: flight mode and ground state. \\
\topic{vehicle_attitude_setpoint} & Output: attitude $\vect q_d$ and collective thrust. \\
\topic{indi_rate_feedforward} & Output: body-rate feed-forward for the rate loop. \\
\topic{vehicle_local_position_setpoint} & Output: setpoints and thrust vector, used by the land detector. \\
\topic{indi_pos_status}, \topic{takeoff_status} & Output: logging and take-off state. \\
\bottomrule
\end{tabularx}
\end{table}

\section{The acceleration law}\label{sec:outer_law}

\subsection{Frames and symbols}

The world frame is north-east-down.
The body frame is forward-right-down.
$R$ is the rotation from body to world.
The gravity vector is $\vect g = (0,0,+g)$.
The vector $\vect z_b$ is the body down axis.
The accelerometer measures the specific force
\begin{equation}\label{eq:outer_fb}
  \vect f_b = R^{\mathsf T}\,(\vect a - \vect g) .
\end{equation}
The measured acceleration is therefore $\vect a_f = R\,\vect f_{b,f} + \vect g$, where $\vect f_{b,f}$ is the filtered specific force of \cref{eq:outer_fb}.

The control variable is the specific thrust vector $\vect T$, in m/s$^2$.
It is the thrust per unit of mass.
A multicopter pushes along $-\vect z_b$, so $\vect T = -\tau\,\vect z_b$ with the thrust magnitude $\tau = \lVert\vect T\rVert > 0$.
At hover, $\tau = g$.

\indidiagram{outer_loop}{Acceleration-loop INDI of the position controller. The stock position/velocity cascade (integrator off) gives a reference acceleration $\aref$. The measured specific force gives $\vect a_f$, and the error is added to the filtered thrust vector $\vect T_f$ to form $\vect T_c$. After saturation the thrust vector sets the attitude setpoint and the collective thrust. A jerk feed-forward drives the rate loop directly. On the ground, in the take-off ramp and with stale data the force source switches to a modelled force.}{fig:outer_loop}

\Cref{fig:outer_loop} shows the loop.
The measured specific force closes the acceleration loop.
The position and velocity cascade closes the outer loop around it.
The blocks marked $H$ are the matched low-pass filters of \cref{sec:outer_filters}.
The force-source switch selects the measured force (flight) or the modelled force (ground, take-off ramp, stale data), as described in \cref{sec:outer_state}.
The thrust-state block forms $\vect T_f$ from the measured and the modelled rotor state with the blend of \cref{sec:outer_gap}.

\subsection{Reference acceleration}\label{sec:outer_cascade}

The reference comes from the stock cascade, with its integrator off.
$K_p$ and $K_v$ are diagonal gain matrices.
The equations show the horizontal entries; the vertical entries are the stock gains.
\begin{align}
  \vect v_{\mathrm{sp}} &= K_p\,(\vect p_{\mathrm{ref}} - \vect p) + \vect v_{\mathrm{ref}} , \label{eq:outer_vsp}\\
  \aref &= K_v\,(\vect v_{\mathrm{sp}} - \vect v) + \vect a_{\mathrm{ff}} . \label{eq:outer_aref}
\end{align}
The velocity setpoint is limited by the stock velocity limits.
The stock filters stay in place.
The integrator is set to zero in every cycle.
The INDI increment of \cref{eq:outer_inc} supplies the integral action.

The horizontal gains $K_p$ and $K_v$ are \param{MC_INDI_XY_P} and \param{MC_INDI_XY_VP}.
They replace the stock horizontal gains, so the user can tune them for the INDI loop.
The vertical gains stay stock.
The feed-forward acceleration $\vect a_{\mathrm{ff}}$ comes from the setpoint after the take-off and failsafe edits.
A setpoint that arrives at a lower rate than the controller is held and extrapolated.
With $\delta$ the age of the setpoint, the velocity grows by $\vect a\,\delta$ and the position by $\vect v\,\delta + \tfrac12 \vect a\,\delta^2$.
$\delta$ is capped at 0.5~s.

\subsection{Increment and thrust inversion}\label{sec:outer_inc}

The filtered thrust vector of the actuators is $\vect T_f$.
It is built from the rotor speed (\cref{sec:outer_state}).
The law adds the acceleration error to it:
\begin{equation}\label{eq:outer_inc}
  \vect T_c = \vect T_f + (\aref - \vect a_f) .
\end{equation}
This is the incremental form of \cref{eq:indi_law}.
The control effectiveness is one, because $\vect T$ is already a specific thrust.
No model of the airframe appears in \cref{eq:outer_inc}.

The actuator signal is a normalised rotor speed squared, $u_i = (\Omega_i/\Omega_{\max})^2$ as in \cref{eq:u_state}.
The thrust model needs a scale, $g_t$ of \cref{eq:outer_gt}.
It is
\begin{equation}\label{eq:outer_gt}
  g_t = g / h_t , \qquad \vect T_f = -g_t\,u_f\,\vect z_b ,
\end{equation}
where $u_f$ is the filtered mean actuator state of all rotors (a scalar; in \cref{ch:inner_loop} $u_f$ is the torque).
$h_t$ is the hover thrust.
In flight, while the hover-thrust estimate is valid, $h_t$ is that estimate.
Otherwise it is \param{MPC_THR_HOVER}.
It is never below 0.05.
An error of $g_t$ changes the loop gain.
It does not leave an offset (\cref{sec:outer_gap}).

The normalised collective $T_{\mathrm{col}}$ of \cref{eq:outer_cmd} that goes to the attitude setpoint is
\begin{equation}\label{eq:outer_cmd}
  T_{\mathrm{col}} = \frac{\tau}{g_t}, \qquad \tau = \lVert\vect T_c\rVert ,
\end{equation}
limited to the range from \param{MPC_THR_MIN} to \param{MPC_THR_MAX}.
Until the take-off state is flight, the lower limit is zero.

\begin{keyresult}[No thrust model in the loop gain]
The model scale $g_t$ only enters through $\vect T_f$ and through the final map to the collective.
A wrong hover thrust therefore changes the loop gain.
It does not leave a standing acceleration error.
\end{keyresult}

\subsection{Filters and delay matching}\label{sec:outer_filters}

The accelerometer and the measured thrust pass through the same second-order low-pass filter.
Its cut-off is \param{MC_INDI_POS_LPF}.
The filter on the force runs at the accelerometer sample rate.
The module is scheduled on the acceleration topic.
The force filter runs on every acceleration sample; the stage itself, the INDI law, runs once per update of the position estimate.
This is on purpose.
Filtering after a decimation to the lower stage rate would fold a vibration line into the signal band.
The stage runs at about 89~Hz in the reference simulation, so a simulated vibration line at 95.7~Hz would appear at 6.8~Hz in $\vect a_f$.
The filter on the thrust runs at the stage rate (about 89~Hz in the reference simulation) with the same cut-off, so both paths have the same magnitude and phase.

The ESC-based thrust arrives later than the accelerometer.
A pure delay of \param{MC_INDI_POS_DLY} controller cycles (0 to 4) matches it to the force path.
The modelled thrust path gets the same delay.

\subsection{Saturation and attitude setpoint}\label{sec:outer_sat}

The thrust vector $\vect T_c$ is limited in this order.
\begin{enumerate}
  \item Upward thrust is mandatory.
    If the vertical part is not clearly upward, it is set to a small upward value.
  \item The tilt is limited to \param{MPC_TILTMAX_AIR} in flight and to \param{MPC_TILTMAX_LND} before flight, with a slew between the two.
    The horizontal part is scaled down to fit.
  \item If $\tau$ is above its maximum, the horizontal part is shrunk first and the vertical part is kept.
    If $\tau$ is below its minimum, the whole vector is scaled.
    In both cases the blend of \cref{sec:outer_gap} is held in the next cycle.
\end{enumerate}
The desired body axis is $\vect z_{b,d} = -\vect T_c/\lVert\vect T_c\rVert$.
The attitude setpoint $\vect q_d$ follows from $\vect z_{b,d}$ and the yaw setpoint.
If the setpoint has no yaw, the current yaw is used.
The yaw rate setpoint is passed on.
The vertical axis keeps priority over the horizontal one: a saturated vehicle holds its altitude before it holds its path.

\subsection{Jerk feed-forward to the rate loop}\label{sec:outer_ff}

The attitude loop follows $\vect q_d$ with a lag.
A change of the acceleration setpoint would reach the rate loop late.
The module therefore sends the body-rate feed-forward of \cref{eq:outer_ff}
\begin{equation}\label{eq:outer_ff}
  \vect\omega_{\mathrm{ff}} = \frac{1}{\tau}\,\bigl(\vect j\cdot\vect y_b,\; -\vect j\cdot\vect x_b,\; 0\bigr) ,
\end{equation}
where $\vect j$ is the jerk and $\vect x_b$, $\vect y_b$ are the desired body axes.
It is limited to 3~rad/s per axis, and it is zero when $\tau$ is below 1~m/s$^2$.
The rate loop adds it to its setpoint (\cref{ch:inner_loop}).

The jerk is the difference of the setpoint acceleration $\vect a_{\mathrm{ff}}$ over time.
The module computes it itself.
It is held for 0.1~s between changes, limited to 50~m/s$^3$ and low-pass filtered with the cut-off of the force filter.
The jerk field of the trajectory setpoint is not used.
In the automatic modes that jerk is the bang-bang jerk of the trajectory generator.
It is not the derivative of the published acceleration.
The feed-forward is active in every mode.
It is published as valid only when the stage runs on the measured force, in flight, not landed, and not in climb-rate-only mode.

\section{Measured and modelled state}\label{sec:outer_state}

The stage runs in one of four modes (\cref{tab:outer_stage}).
The mode selects where the force and the thrust come from.

\begin{table}[tbp]
\caption{Operating modes of the INDI stage.}\label{tab:outer_stage}
\centering
\small
\begin{tabularx}{\textwidth}{@{}L{2.4cm} L{5.6cm} Y@{}}
\toprule
Mode & Condition & Law input \\
\midrule
Bypass & INDI not requested: disarmed, or no position, velocity or climb-rate mode. & The stage does not run. Its state follows the commanded collective. \\
Hold & Requested, but the attitude is stale, $\aref$ is not finite, the cycle interval is above 0.2~s, or the stage rejected its input. & The last INDI attitude and thrust are held. \\
Modelled force & Requested, acceleration data stale, or ground and take-off ramp. & Modelled specific force. \\
Measured & Attitude and acceleration fresh, in flight, no ground contact. & Measured specific force and measured rotor speed. \\
\bottomrule
\end{tabularx}
\end{table}

\paragraph{Measured thrust.}
The mean actuator state is $\bar u = \frac14\sum_i (\Omega_i/\Omega_{\max})^2$ over the four rotors.
$\Omega_i$ comes from the ESC rotor-speed telemetry, and $\Omega_{\max}$ is \param{MC_INDI_W_MAX}.
\param{MC_INDI_W_MAX} must be set: with its default of 0 the INDI rate module refuses to arm, and the outer loop has no measured thrust.
The mean is for four rotors (a quadrotor).
If a motor goes offline in flight, or the ESC data is stale, the value is invalid.
The law then uses the commanded collective instead.

\paragraph{Modelled force.}
When the measured force cannot be used, the force is computed from the delayed actuator signal itself, with the same filter.
The measured and modelled paths are then the same signal.
They cancel in \cref{eq:outer_inc}, and the thrust becomes $\vect T_c = \aref - \vect g$.
The stage then acts as a feed-forward from acceleration to thrust.
This is the safe source on the ground, where the accelerometer also holds the ground reaction.

\paragraph{Freshness.}
Each signal has an age limit, and a stale signal acts in its own way.
\begin{itemize}
  \item Attitude or acceleration older than 20~ms changes the mode (\cref{tab:outer_stage}).
  \item ESC data older than 100~ms makes only the measured thrust invalid. The law then runs on the modelled thrust and the blend is frozen. The force still comes from the accelerometer.
  \item Allocator status older than 200~ms is ignored.
  \item The age of the commanded collective only decides which value restarts the stage.
\end{itemize}
After an acceleration loss, the measured path returns only after 20~ms of continuously fresh acceleration.

\paragraph{Continuity at switches.}
While the telemetry is invalid, the measured thrust path is fed the modelled value minus the blend, so the thrust stays continuous.
When the force becomes valid again, its filter restarts from the force used in the last cycle.
It then converges at the filter bandwidth.
A reset of the stage re-seeds the delay line and the thrust filter from the measured rotor state when it is valid, and else from the commanded collective minus the blend.

\section{Removing a steady thrust gap}\label{sec:outer_gap}

If a motor weakens or the battery sags, the rotor speed no longer gives the thrust that the model expects.
The outer loop has no integrator, so a standing acceleration error would remain.
The module restores integral action with a complementary blend of measured and modelled collective.
It adds a slow correction $b$ to the measured actuator state (\cref{eq:outer_blend}):
\begin{equation}\label{eq:outer_blend}
  b \leftarrow b + \frac{\ts}{t_b + \ts}\,\bigl(u_{\mathrm{mod},f} - u_{\mathrm{meas},f} - b\bigr),
  \qquad u_f = u_{\mathrm{meas},f} + b .
\end{equation}
Here $\ts$ is the stage period and $u_{\mathrm{meas},f}$ is the filtered mean of the four rotor states (a scalar, unlike the torque vector of \cref{ch:inner_loop}).
$u_{\mathrm{mod},f}$ is the filtered commanded collective, which the thrust model predicts (the model path of the thrust).
The blend is bounded to $\lvert b\rvert \le 0.5$.
The time constant is $t_b = $ \param{MC_INDI_POS_TB}.

At high frequency the measured rotor speed dominates $u_f$.
The loop reacts to fast thrust changes without the model.
At low frequency $u_f$ equals the model.
With a thrust gap, the measured rotor state is only a fraction $\rho$ of the commanded state.
Without the blend, the loop of \cref{eq:outer_inc} would settle at $\vect a_f = \aref - (1-\rho)\,\vect T_c$.
The blend brings $u_f$ back to the model at low frequency and removes this error.
The result is an integral action with the time constant $t_b$.
A wrong scale $g_t$ is different: it changes only the loop gain, so it leaves no offset.
The bound of 0.5 is large enough to absorb the thrust drop that follows a sudden loss of hover thrust.

The blend is frozen in these cases:
\begin{itemize}
  \item the vehicle is landed or maybe landed, or the stage uses the modelled force;
  \item the telemetry is invalid;
  \item the collective was clamped in the previous cycle (this prevents windup);
  \item the measured mean actuator state is at or above 0.98;
  \item the allocator reports saturation.
\end{itemize}
A reset of the stage keeps the blend, because it is a property of the airframe and battery.
Only arming and disarming clear it.

\begin{keyresult}[Integral action without an integrator state in the cascade]
The cascade has no integrator.
The blend gives the outer loop the low-frequency behaviour of one, and its freeze rules give it an anti-windup.
\end{keyresult}

\section{Setpoint sources and flight modes}\label{sec:outer_modes}

The module keys on the control-mode flags, not on mode names.
Every mode in which the stock position controller would run goes through the same stage (\cref{tab:outer_modes}).
The setpoints come from the unchanged flight-mode manager and from offboard sources.
Both publish \topic{trajectory_setpoint}.

\begin{table}[tbp]
\caption{Control-mode flags and use of the INDI stage.}\label{tab:outer_modes}
\centering
\small
\begin{tabularx}{\textwidth}{@{}L{5.4cm} Y@{}}
\toprule
Flags set & Use of the stage \\
\midrule
Position or velocity control: Position, Hold, Mission, RTL, Land, Take-off, Orbit, Follow, offboard position, velocity or acceleration, goto & Full INDI. $\aref$ comes from the cascade and the thrust direction from the INDI stage. \\
Climb-rate control only: Altitude, Descend failsafe & INDI collective only. The direction is that of the stock thrust vector, and the rate feed-forward is zero. The stock direction must be finite and at least $5.7^\circ$ above the horizon, else the stage holds. \\
None of these: Manual, Stabilized, Acro, offboard attitude or rate & Not served by this module. The stock attitude and rate paths act. \\
\bottomrule
\end{tabularx}
\end{table}

The setpoint sources are the following.
\begin{itemize}
  \item Position, Hold, Mission, RTL, Land and Take-off: the flight-mode manager gives position, velocity and acceleration references.
  \item Goto and offboard: the setpoint comes from the sender, held and extrapolated between samples.
  \item Climb-rate modes: the vertical velocity setpoint only.
\end{itemize}
In every case the cascade of \cref{eq:outer_vsp,eq:outer_aref} turns the setpoint into $\aref$.
The INDI stage does not know the mode.

\section{Ground, take-off and landing}\label{sec:outer_ground}

One law runs from arming to disarming.
The stock thrust is never published.
The stage is requested when the vehicle is armed and a position, velocity or climb-rate mode is active.
\Cref{fig:outer_ground} shows the states.

\indidiagram{ground_logic}{Operating states of the outer loop. The same INDI law runs from arming to disarming; only the source of the force and thrust state changes. The note at the edge applies to the step discussed under ``Edge to flight''. Ground and ramp use the modelled force, flight uses the measured force and rotor speed, and touch-down returns to the modelled force and holds the blend.}{fig:outer_ground}

\paragraph{On the ground.}
The accelerometer holds the ground reaction, so the measured force would be wrong.
The stage uses the modelled force.
The module replaces the setpoint by an empty one with a large downward acceleration of 100~m/s$^2$, also in flight while ground contact is reported.
This drives the thrust to the minimum.
The tilt limit is \param{MPC_TILTMAX_LND}, and the horizontal acceleration is decoupled.

\paragraph{Take-off ramp.}
The stock ramp (\param{MPC_TKO_RAMP_T}, \param{MPC_TKO_SPEED}) shapes the vertical speed.
The collective rises from zero through the model inversion.
The vertical feed-forward acceleration is cleared during the ramp when a vertical velocity setpoint is given.
On the ground and in the ramp these parts are frozen: the blend, the rate feed-forward and the stock integrator.

\paragraph{Edge to flight.}
At the end of the ramp the force source switches from modelled to measured.
The step of the thrust at this edge is the error of the thrust model, a few percent of the hover thrust.
In the take-off after a re-arm the step was 4 to 5~\% of the hover thrust; every take-off stayed below a limit of 10~\%.

\paragraph{Landing.}
The stock land task and \param{MPC_LAND_SPEED} shape the setpoint.
At touch-down the landed flags switch the stage to the modelled force and hold the blend.
The ground reaction then cannot wind anything up.
While descending or climbing at a commanded vertical speed without altitude control, the vertical measurement uses the altitude derivative, weighted by the ratio of the vertical speed setpoint to \param{MPC_LAND_SPEED}.

\paragraph{Arm and disarm.}
Both edges restart the stage state: blend, held setpoint, stale latch and jerk differencing.

\section{Estimator resets and failure signals}\label{sec:outer_resets}

\paragraph{Estimator resets.}
When the estimator resets its position, velocity or heading, the stock logic corrects the setpoints and the position-controller filters.
The module clears only the jerk history.
It keeps the measured force, the mean actuator state, the delay lines and the previous thrust.
These signals do not depend on the world frame.
In the reference simulation a single raw accelerometer sample carries vibration of about 0.25~m/s$^2$, against 0.04~m/s$^2$ after the filter.
Re-seeding from it would step the acceleration estimate and tilt the thrust command (about 2~deg in simulation).
A full reset of the stage happens only at entry to position control, at the return from hold and when the stage restarts.
It keeps the blend and re-seeds the thrust path as described in \cref{sec:outer_state}.

\paragraph{Failure signals.}
If the setpoint is stale or the cascade fails, the module first tries a fallback within 200~ms of the last valid setpoint.
After that it uses a failsafe setpoint.
\begin{itemize}
  \item Horizontal velocity known: hold zero horizontal velocity, the stop-and-wait case.
    Not known: zero acceleration, and descend at the land speed, the blind-land case.
  \item Vertical velocity known: hold zero vertical velocity, except in blind land.
    Not known: a downward acceleration of 0.3~m/s$^2$, a slow descent just below hover.
\end{itemize}
The failsafe setpoint goes through the same INDI stage.
Invalid attitude, acceleration or $\aref$ lead to the hold mode of \cref{tab:outer_stage}.
They never lead to the stock thrust.

\section{Guards}\label{sec:outer_guards}

The module has these protections.
\begin{itemize}
  \item Stage hold. Stale or invalid inputs freeze the last attitude and thrust. A stage interval above 0.2~s holds the stage; shorter intervals are clamped to 2~ms.
  \item Saturation order. Altitude has priority over the horizontal path (\cref{sec:outer_sat}).
  \item Blend freeze. The blend is frozen on clamps, allocator saturation, telemetry loss, high rotor speed and on the ground (\cref{sec:outer_gap}).
  \item Blend bound. $\lvert b\rvert \le 0.5$.
  \item Rate feed-forward limits. 3~rad/s per axis, jerk 50~m/s$^3$.
  \item Measured-path recovery. 20~ms of fresh acceleration.
  \item Missing rotor speed. If a motor goes offline or the ESC data is stale, the stage runs on the modelled collective.
\end{itemize}

\section{Parameters}\label{sec:outer_params}

\Cref{tab:outer_params} lists the parameters of the outer-loop module.
The ``Reference'' column is the set of the simulated reference airframe.
Where that set does not define a parameter, the reference value is the firmware default.

\begin{table}[tbp]
\caption{Parameters of the outer-loop module.}\label{tab:outer_params}
\centering
\small
\begin{tabularx}{\textwidth}{@{}L{3.5cm} r r l Y@{}}
\toprule
Parameter & Reference & Default & Unit & Meaning \\
\midrule
\param{MC_INDI_POS_EN} & 1 & 0 & -- & Start the INDI outer loop instead of the stock one (reboot). \\
\param{MC_INDI_POS_LPF} & 10.0 & 10.0 & Hz & Cut-off of the identical filters on force and measured thrust. 0 turns both off. \\
\param{MC_INDI_POS_DLY} & 1 & 1 & cycles & Delay of the measured thrust (0 to 4). \\
\param{MC_INDI_POS_TB} & 0.3 & 0.3 & s & Blend time constant $t_b$. 0 turns the blend off. \\
\param{MC_INDI_XY_P} & 1.2 & 1.2 & 1/s & Horizontal position gain $K_p$, instead of \param{MPC_XY_P}. \\
\param{MC_INDI_XY_VP} & 2.4 & 2.4 & 1/s & Horizontal velocity gain $K_v$, instead of \param{MPC_XY_VEL_P_ACC}. \\
\param{MC_INDI_W_MAX} & 2700 & 0 & rad/s & Rotor speed at full thrust $\Omax$ (shared with the rate module). Must be set; the default 0 blocks arming of the INDI rate loop. \\
\bottomrule
\end{tabularx}
\end{table}

Eight more parameters (\param{MC_INDI_PINJ_*}) drive a bench excitation for loop-break measurements.
They are off by default and are not part of the flight design.
The module also uses stock parameters, without changing them.
These are the hover thrust \param{MPC_THR_HOVER}, the thrust limits \param{MPC_THR_MIN} and \param{MPC_THR_MAX}, the vertical gains, the tilt limits \param{MPC_TILTMAX_AIR} and \param{MPC_TILTMAX_LND}, the velocity limits, and the take-off and landing parameters.
On the reference airframe \param{MPC_THR_HOVER} is 0.137 and \param{MPC_THR_MIN} is 0.05.
The stock defaults of the tilt limits are 45 and 12~deg.
The ramp defaults are 3.0~s and 1.5~m/s.

\paragraph{Library constants.}
Some values are fixed in the firmware: a blend bound of 0.5, a rotor-speed limit of 0.98 for the blend, a jerk limit of 50~m/s$^3$, a rate feed-forward limit of 3~rad/s, a jerk hold of 0.1~s, at most 4 delay samples, and a minimum upward thrust of 0.1~m/s$^2$, and no rate feed-forward below a thrust of 1~m/s$^2$.

\section{Outer-loop margins}\label{sec:outer_margins}

The stability margins of the position and velocity loop around the acceleration loop are measured and discussed in \cref{ch:analysis}.
They cover the horizontal axes and the vertical axis, with the gains, filters, delay and blend of \cref{tab:outer_params}.
